%% sections/05_discussion.tex

\section{Discussion}
\label{sec:discussion}

\subsection{Efficiency Gain for High-Eccentricity Orbits}
\label{sec:efficiency}

The $r^{3/2}$ step scaling (Equation~\ref{eq:dt_keplerian}) implies
that \AAS\ concentrates computational effort at perihelion where
orbital curvature is highest. For Phaethon with $e = 0.890$, the
perihelion and aphelion distances are $q = 0.140$\,AU and
$Q = 2.402$\,AU respectively, giving a step ratio
$\Delta t_Q / \Delta t_q \approx (Q/q)^{3/2} \approx 117$.
A fixed-step integrator must use the perihelion step everywhere,
wasting a factor of $\sim 100$ function evaluations at aphelion.
To avoid ambiguity, we distinguish three non-equivalent efficiency
definitions: (a) cost at fixed accuracy, (b) accuracy at fixed cost,
and (c) wall-clock time on a specified machine.
In the present dataset we report (a) using function-evaluation counts.
For Phaethon at $|\Delta E/E|\approx7.00315\times10^{-5}$,
\AAS\ requires $7.0644\times10^4$ evaluations
($\varepsilon=10^{-3}$) versus $2.922\times10^4$ for \saba,
i.e.
a factor $2.42$ higher cost at equal accuracy.
Accordingly, this benchmark should not be interpreted as a global
speedup claim for \AAS; rather, it quantifies a trade-off between
geometric fidelity and computational cost in the tested regime.
The efficiency advantage of adaptive step-size symplectic integration
for high-eccentricity orbits is also documented in recent works
\citep{Wang2024tts, Ye2025}.

\subsection{Comparison with \orbitn\ and \rebound}
\label{sec:comparison}

\citet{Zeebe2023} showed that \orbitn\ achieves
$|\Delta E/E| \sim 10^{-11}$--$10^{-12}$ for solar-system
integrations over 100\,Myr with Kahan compensated summation.
Our benchmark set confirms that \AAS\ remains in the
$10^{-13}$ regime on long arcs (Apophis, 50\,yr:
$1.96\times10^{-13}$), while fixed-step \saba\ degrades for
high-eccentricity trajectories.

The \saba\ family \citep{LaskarRobutel2001, Farres2013}
provides excellent accuracy for nearly circular orbits at fixed step,
and can improve by up to six orders of magnitude over standard WH
\citep{ReinHernandezTamayo2019}. \AAS\ is not intended to supersede
\saba\ for circular-orbit problems; its advantage emerges for
high-eccentricity objects where adaptive stepping recovers efficiency.

The 1PN corrections in \AAS\ follow \citet{SahaTremaine1994},
consistent with \orbitn. A detailed validation of the full \astdyn\
force model (N-body, 1PN, SRP, Yarkovsky) against JPL Horizons is
reported in the companion paper (Bigi, in prep.).

\paragraph{Comparison with IAS15.}
We include a direct comparison with \textsc{ias15}
\citep{ReinSpiegel2015}, the 15th-order Gauss--Radau
integrator implemented in \textsc{rebound} \citep{ReinLiu2012}.
The Taylor-series-based \textsc{heyoka} integrator
\citep{Biscani2021heyoka} represents a further non-symplectic
alternative achieving sub-machine-precision accuracy for
astrodynamics problems.
\textsc{ias15} is the standard high-accuracy reference for
single-body integrations; it is not symplectic and does not
preserve a shadow Hamiltonian.

For a meaningful comparison, we adopt a common diagnostic:
the relative variation of the two-body Keplerian energy
$|\Delta H_{\rm 2b}/H_{\rm 2b,0}|$, which measures the
pure numerical error of each integrator on the unperturbed
Keplerian problem. On a 50-year integration of Apophis,
\textsc{ias15} achieves $|\Delta H/H_0| = 1.69\times10^{-16}$
(near machine precision), compared with
$1.96\times10^{-13}$ for AAS,
$9.46\times10^{-8}$ for SABA4, and
$6.74\times10^{-14}$ for \textsc{rkf7(8)}.
These values reflect the characteristic behaviour of each
class: \textsc{ias15} minimises local truncation error
by construction; AAS preserves the shadow Hamiltonian
through symplecticity; SABA4 accumulates fixed-step
splitting error for high-eccentricity orbits; and
\textsc{rkf7(8)} reaches near-machine accuracy on short
arcs through adaptive tolerance control.

The time-reversal test (two-body, 50~yr) gives
$\epsilon_r = 4.22\times10^{-11}$ for AAS,
$3.95\times10^{-12}$ for SABA4,
$2.98\times10^{-11}$ for \textsc{rkf7(8)}, and
$4.53\times10^{-13}$ for \textsc{ias15}.
All four integrators achieve sub-$10^{-10}$ reversibility
on the unperturbed Keplerian problem, confirming that
the time-reversal diagnostic on the two-body system does
not distinguish symplectic from non-symplectic integrators
at this precision level. The distinction emerges in the
long-term energy behaviour: AAS maintains
$|\Delta H/H_0| < 10^{-12}$ across all 12 benchmark objects
and all integration arcs (50--10\,000~yr), while non-symplectic
integrators show secular accumulation in the perturbed
N-body regime.

The maximum Lyapunov exponent (mLCE) shows agreement
within 8\% across all four integrators for Apophis
($\lambda \approx 0.13$~yr$^{-1}$), confirming that
none of the integrators introduces spurious numerical chaos.

\subsection{Kahan Compensated Summation}
\label{sec:kahan}

Following \citet{Kahan1965} and \citet{Zeebe2023}, \AAS\ implements
Kahan compensated summation in the force accumulation loop.
A dedicated with/without-Kahan ablation is not included in the
current benchmark CSV set; we therefore defer a quantitative
isolation of Kahan's contribution to future work.
Consequently, no performance or accuracy claim in this manuscript
should be interpreted as being individually attributed to Kahan
summation.

\subsection{Statistical Robustness and Sensitivity}
\label{sec:stats_sensitivity}

The present benchmarks are deterministic trajectory comparisons and
do not yet provide full statistical uncertainty quantification.
In particular, confidence intervals, distributional robustness across
initial-condition perturbation ensembles, and parameter-sensitivity
maps (with respect to $\varepsilon$, force-model toggles, and
finite-difference scales in STM validation) are not yet reported.
These analyses are required to convert the current point estimates
into statistically characterized performance claims.

\subsection{Methodological Scope vs Software Scope}
\label{sec:method_scope}

The short-term benchmark suite (Section~\ref{sec:benchmark}) uses
intentionally reduced-complexity two-body diagnostics: they isolate
numerical behavior (energy drift, reversibility, STM consistency)
under controlled conditions, but are not by themselves sufficient
to establish operational reliability in strongly perturbed
multi-body regimes. This choice enables clean mechanism-level
validation; the limitations are discussed in
Section~\ref{sec:limitations}.

Because \astdyn\ is a full pipeline library, manuscript narrative can
easily drift toward implementation details. In this paper, our intent
is methodological: geometric properties of the integrator,
error-behavior diagnostics, and reproducible benchmark definitions.
Software-engineering and pipeline-deployment details are treated as
supporting context and are delegated, where possible, to the
companion library paper.

\subsection{Application to Occultation Prediction}
\label{sec:occultation}

For short-arc occultation workflows, the 30-day Horizons residuals
from \texttt{horizons\_short.csv} remain at
$10^{-6}$--$10^{-4}$ AU scale in the present benchmark setup, e.g.
for Apophis $\delta r(30\,\mathrm{d})=2.92913\times10^{-6}$ AU
(\AAS), $2.25895\times10^{-4}$ AU (\saba), and
$2.92913\times10^{-6}$ AU (\rkf).
These are external-reference ephemeris offsets and should not be read
as pure integrator local error; they also include modeling/setup
differences relative to Horizons.
For uncertainty propagation, the 30-day comparison in
\texttt{uncertainty.csv} gives
$\sigma_{AT}=9.25761\times10^{-9}$ AU (CovProp) and
$9.48554\times10^{-9}$ AU (Monte Carlo), with
$97.60\%$ agreement, supporting the validity of the STM-based
linearized covariance model over this time window.

\subsection{Long-Term Dynamical Fidelity}
\label{sec:lt_discussion}

The long-term test suite of Section~\ref{sec:longterm}
establishes three properties that are essential for
operational use of \AAS\ in occultation prediction pipelines.

\textit{Absence of secular drift.} The linear regression
analysis (Section~\ref{sec:lt_secular}) confirms that
\AAS\ produces bounded, oscillatory energy errors rather
than linear secular accumulation. This behavior is consistent
with geometric integration in the symplectic-kernel framework
\citep{HairerLubichWanner2006},
ensures that propagation errors do not grow monotonically
with arc length --- a critical requirement for multi-year
orbital solutions.

\textit{Absence of spurious chaos.} The agreement between
\AAS\ and \rkf\ on the mLCE values
(Section~\ref{sec:lt_lyapunov}) confirms that \AAS\ does
not introduce numerical chaos beyond what is physically
present in the orbit. The Griqua result is particularly
informative: both integrators agree on $\lambda \sim
10^{-3}$\,yr$^{-1}$, consistent with the known chaotic
timescale of the $2:1$ resonance \citep{ferrazmello1994}.

More generally, the methodological contribution of \AAS\ is
best interpreted as a reformulation and integration of known
adaptive-geometric principles (rather than a new adaptive
scaling law): the novelty lies in the Hessian-spectral-radius
formulation, its curvature-based interpretation, and its direct
reuse for analytical STM propagation.

\textit{Geometric conservation diagnostics.} The Jacobi constant
test (Section~\ref{sec:lt_jacobi}) and the time-reversal
test (Section~\ref{sec:lt_reversibility}) directly probe
the geometric structure of the integration. \saba\ performs
similarly to \AAS\ on these diagnostics, while \rkf\ shows the systematic
irreversibility characteristic of adaptive Runge-Kutta
methods.

Importantly, these diagnostics should not be read as a standalone
proof of symplecticity in physical time: as discussed in
Section~\ref{sec:symplecticity_proof}, time-reversibility and bounded
energy oscillations are necessary but not sufficient to claim
strict symplecticity of the adaptive map.

The practical implication for \astdyn\ users is that
\AAS\ can be used with confidence for arcs up to
50\,yr for NEA/PHA objects and up to
1000\,yr for Trojans and 10000\,yr for TNOs, with energy errors
remaining below $10^{-11}$ throughout.

\subsection{Limitations and Future Work}
\label{sec:limitations}

The current \AAS\ implementation has the following limitations.

\textit{Close encounters.} The Hessian metric becomes inaccurate
when inter-body forces are comparable to the central-body force.
A switching strategy such as that of \citet{HernandezDehnen2024}
or time-reversible close-encounter regularization \citep{Lu2024trace}
would be required for asteroid--planet flyby scenarios.

\textit{N-body perturbations.} When planetary perturbations are
active, $\rhoH$ accounts only for the central-body Hessian.
Incorporating the full force Jacobian would improve the step estimate
in strongly perturbed regions (e.g., Jupiter resonance crossings).

\textit{Operational exploitation of the analytical STM.} The analytical
STM machinery is implemented and internally validated, but its full
practical leverage in orbit determination has not yet been exhausted
in this manuscript (e.g., covariance realism tests, estimator
conditioning studies, and residual-level comparisons against numerical
STM pipelines).

\textit{Benchmark representativeness.} A substantial fraction of the
current validation isolates two-body dynamics to separate integrator
effects from force-model complexity. While this is useful for
mechanism-level diagnostics, it is not fully representative of
operational ephemeris propagation. A dedicated benchmark extension is
therefore required with full N-body dynamics (at minimum including
Jupiter), strong-perturbation intervals (resonance crossings), and
close-encounter scenarios where adaptive behavior is most stressed.

\textit{SABA-type correctors.} \citet{LaskarRobutel2001} showed
that symplectic correctors can improve accuracy by orders of
magnitude at no integration cost. Combining \AAS\ with the SABAC
corrector step is a natural extension.

\textit{Clarify the symplectic nature of the method.} Because
\AAS\ uses an adaptive step function, the exact mathematical
scope of the symplectic claim should be stated more explicitly
(e.g., conditions under which the map is strictly symplectic,
approximately symplectic, or only time-reversible). If needed,
the claim should be weakened in the manuscript wording to match
the proven properties.

Overall, the current status of the method can be summarized as
technically solid and scientifically promising, but not yet fully
mature in the sense of complete theory-to-operations validation.
